% !TEX root = ../main.tex

\section{Identification and Position}
\label{sec:results-identification}

With a working signal chain, every pedal at hand was identified with the correct class and wiring (\cref{tab:pedals}). The M-Audio EX-P carries 8 to \qty{15}{\percent} of its total resistance in its wiper lead. With the original wiper limit of \qty{10}{\percent}, it was therefore classified as other over most of its travel, which led to the limit of \qty{25}{\percent} and the near-end class. A bare \qty{100}{\kilo\ohm} potentiometer between tip and sleeve stood in for a rheostat pedal, which was not available, and several sustain pedals behaved like the Boss FS-5U footswitch. A pedal with the wiper on the tip, connected through a mono cable, was classified as a rheostat or as other and jumped between minimum and maximum when moved, as expected from its parallel track halves (\cref{sec:network-model}). No systematic plug-cycle test was made. In use, occasional misdetections occurred, most often a footswitch taken for a rheostat when it was plugged in slowly or turned in the jack, and the firmware recovered on its own within a few seconds.

\begin{table}[tb]
  \centering
  \footnotesize
  \titledcaption{Identified Pedals}{Classification of the available pedals. Solver total: total resistance from the solver, including the wiper lead, against the track measured with an ohmmeter. For the footswitch, the closed resistance is compared against the threshold.}
  \label{tab:pedals}
  \begin{tabularx}{\linewidth}{@{}>{\raggedright\arraybackslash}X l l r r@{}}
    \toprule
    \textbf{Pedal} & \textbf{Cable} & \textbf{Class} & \textbf{Solver total} & \textbf{Ohmmeter} \\
    \midrule
    Boss EV-30, wiper on tip & stereo & potentiometer & \qty{10.85}{\kilo\ohm} & \qty{9.81}{\kilo\ohm} \\
    M-Audio EX-P, wiper on ring & stereo & potentiometer (near end) & \qty{11.40}{\kilo\ohm} & \qty{11.39}{\kilo\ohm} \\
    Boss FS-5U, closed & mono & footswitch & $\leq \qty{0.2}{\kilo\ohm}$ & \qty{8.7}{\ohm} \\
    \qty{100}{\kilo\ohm} potentiometer, tip to sleeve & mono & rheostat & \qty{100.15}{\kilo\ohm} & \qty{99.98}{\kilo\ohm} \\
    \bottomrule
  \end{tabularx}
\end{table}

The total resistance only sets the settling time and the scale of a rheostat, whereas the position depends on voltage ratios alone. For the EX-P and the rheostat, the solver total agrees with the ohmmeter within \qty{0.1}{\percent} and \qty{0.2}{\percent}. The EV-30 was solved \qty{11}{\percent} high at bring-up, still at \qty{315}{\hertz}, and at about \qty{9.4}{\kilo\ohm} from the voltage levels of a later oscilloscope capture. The EV-30 has a knob that sets the minimum of its range and adds resistance to its network, and this knob was set differently between the measurements, which may explain the difference.

With the EV-30 held still at \qty{819}{\hertz}, the solved position had a standard deviation of \num{274}~ppm of the travel over repeated solves, and a single wiper reading one of \num{203}~ppm (\cref{tab:adc-characterization}). One MIDI step is \num{7874}~ppm, so the spread stays below 1/28 of a step, and with the hysteresis of 0.6 steps, the sent value does not change.
